\documentclass[]{scrartcl}

%-------------------------------------------------------------------------------------------------
% You have to use this package to load the template
\usepackage{Model_Article}

%%%%%%%%%%%%%%%%%%%%%%%%%%%%%%%%%%%%%%%%%%%%%%%%%%%%%%%%%%%%%%%%%%%%%%%%%%%%%%%
% for including lines of code

\usepackage{listings}

%defining colors
\definecolor{string}{RGB}{121, 56, 25}
\definecolor{comment}{RGB}{68, 107, 61}
\definecolor{number}{rgb}{0.6,0.60,0.6}
\definecolor{keyword}{RGB}{4,90,141}
\definecolor{class}{RGB}{4, 123, 93}
\definecolor{method}{RGB}{80, 80, 80}
\definecolor{param}{RGB}{4,90,141}

\lstdefinestyle{mystyle}{
    backgroundcolor=\color{white},   
    commentstyle=\color{comment},
    keywordstyle=\color{keyword},
    numberstyle=\tiny\color{number},
    stringstyle=\color{string},
    basicstyle=\rmfamily\selectfont\normalsize,
    breakatwhitespace=false,         
    breaklines=true,                 
    captionpos=b,                    
    keepspaces=true,                 
    numbers=left,                    
    numbersep=5pt,                  
    showspaces=false,                
    showstringspaces=false,
    showtabs=false,%,tab=\rightarrowfill                  
    tabsize=2, % <-- change tab size here
    emph= {MyClass}, % <-- list your classes here
    emphstyle= \color{class},
    emph= {[2]my_method}, % <-- list your methods here (don't remove the [2])
    emphstyle={[2]\color{method}},
    emph= {[3]my_var},  % <-- list your variables and attributes here (don't remove the [3])
    emphstyle={[3]\color{param}}
}

\lstset{language=Python, style=mystyle} % <-- set language here


% you can define new keywords and commands like that
\newcommand{\myCommand}{\textbf{SpecialCommand}\xspace}

%-------------------------------------------------------------------------------------------------
% here starts actual informations about the content
\title{Minimal Example Template Paper}
\author[1]{G. Laurent}
\affil[1]{Univ. Orléans, CNRS, BRGM, ISTO, UMR 7327, F-45071, Orléans, France}
%\affil[2]{GeoRessources (UMR 7359, Universit\'e de Lorraine / CNRS / CREGU), Vandoeuvre-l\`es-Nancy, F-54518 France}

\Abstract{
	Lorem ipsum dolor sit amet, consectetur adipiscing elit. Sed malesuada vehicula ex, nec consectetur elit malesuada id. Suspendisse ut sagittis eros, ac suscipit velit. In hac habitasse platea dictumst. Curabitur ipsum enim, semper eget sem in, auctor mollis purus. Vivamus viverra aliquet elit, ac iaculis nisl imperdiet aliquet. Pellentesque id ullamcorper eros. Integer vitae enim et sapien fringilla eleifend eu venenatis mauris. Donec placerat euismod lorem, eget semper purus malesuada vitae. Integer fringilla maximus porttitor. Sed euismod ultricies odio in dapibus. Donec consectetur quam eget posuere auctor.
}

\MyKeywords{
	Template \sep Paper \sep Lorem ipsum 
}

% licensing, redefine \lic to change licence, cf. Creative Commons
% by, by-nc, by-nc-nd, or by-nc-sa
\def\lic{by} % cc by is recommended as nc can be exploited by owners of rights to use commercially (publishers)

% status, to indicate the status of the manuscript (draft, preprint, submitted, accepted)
\def\status{accepted}
% this can be change to indicate the target or publishing journal depending on the status
\Journal{Some Journal}
\Volume{96}
\Number{1}
\Pages{145}{169}
\Year{2023}
\DOI{10.1023/j.test.2023}

%-------------------------------------------------------------------------------------------------
% here starts the document, but this is the automated part
\begin{document}

% cover page
\makeatletter
\twocolumn[
\@maketitle
\section*{Highlights}
\begin{compactitem}
\item Highlights are great to make your paper more visible.
\item They should be short sentences.
\item They should give important points.
\end{compactitem}
\vspace{1em}
]
\makeatother

\begin{minipage}{\textwidth}
\section*{Graphical Abstract}
\begin{center}
\includegraphics[width=16cm]{graphical_abstract}
\end{center}

\begin{multicols}{2}
\tableofcontents
%\columnbreak
\end{multicols}
\end{minipage}


\newpage
% end of cover page

\maketitle	
%-------------------------------------------------------------------------------------------------
% start writing from here ↓

\section*{Introduction}

Just put your introduction here.
This is acting as a documentation.
First section illustrates the inclusion of figures (Section~\ref{s:figure}).
There is a section about bibliography (Section~\ref{s:biblio})
 and subsections for backreferencing  (Section~\ref{ss:backref}) and doi integration (Section~\ref{ss:doi}).
Another section is showing how to include programming code (Section~\ref{s:code}).


%sometimes you have to put your figures earlier to make sure they are at the right place
%single column figure (without the *)
\begin{figure}[!b]
	\centering
	\includegraphics[width=9cm]{fig_section} 
	\caption{In sed rutrum felis. Cras convallis libero vel arcu euismod, ut ultricies dolor pellentesque. Nullam congue, mauris et volutpat imperdiet (modified from \citet{Laurent2023}).}
	\label{fig:Figure_col}
\end{figure}

Here is some latin jibberjabber to add some text: 	
	Maecenas volutpat nulla at tellus commodo fermentum. Vestibulum non lacus et lectus consectetur blandit ac eget diam. Sed cursus eros vel mauris consectetur rhoncus non non velit. Vestibulum porta convallis mauris quis varius. Suspendisse eu odio non lacus imperdiet sollicitudin et ut orci. Cum sociis natoque penatibus et magnis dis parturient montes, nascetur ridiculus mus. Vivamus volutpat feugiat congue. Sed scelerisque nulla quis risus gravida mollis. Pellentesque suscipit varius est vitae eleifend. Nam euismod nisi nec diam gravida tristique. Pellentesque habitant morbi tristique senectus et netus et malesuada fames ac turpis egestas. Aenean convallis imperdiet ligula, sit amet finibus lectus maximus eget. Quisque non euismod ex, interdum dignissim tortor. Morbi facilisis metus dolor, eget rhoncus justo hendrerit sed.
	
\section{Figures}
\label{s:figure}

This is to show how to include a figure.
There are single column figures (Fig.~\ref{fig:Figure_col}) and full width figures (Fig.~\ref{fig:Figure_1}).

%full width figure (figure*)
\begin{figure*}[!tbp]
	\centering
	\includegraphics[width=\linewidth]{fig_section} 
	\caption{In sed rutrum felis. Cras convallis libero vel arcu euismod, ut ultricies dolor pellentesque. Nullam congue, mauris et volutpat imperdiet (modified from \citet{Laurent2023}).}
	\label{fig:Figure_1}
\end{figure*}

\section{Bibliography}
\label{s:biblio}

\subsection{Special references backref}
\label{ss:backref}
A back reference is automatically added in the bibliography listing, but it requires an extra compilation to get the page right.

\subsection{DOI and HAL integration}
\label{ss:doi}

There is a command in the template that automatically format the doi references in the bibliography.
All it should take is to make sure the doi is listed in the bib file under the ``doi'' keyword.
It will appear between chevron brackets as a link to doi.org.

If you have listed papers that are in HAL archives you can add them as an extra keyword in the bib file
\begin{lstlisting}
hal = {insu-04182730}
\end{lstlisting}


\section{Code}
\label{s:code}

You can include code inline like this \lstinline!Myclass!. 
For this, the words need to be added to the lstlisting settings at the start.
They can be coloured differently depending on there natures: class, method, attribute or variables.
Comments also have a special colour and strings too, which generally require to have " characters around like in
\lstinline!"A string"!.

Off course, you can also make code listings with the ``lstlisting'' environment:
\begin{lstlisting}[language=Python]
from package import MyClass
a = MyClass()
a.my_var = a.my_method()
print("Example done !")
\end{lstlisting}


%------------------------------------------------------
\begin{appendices}

\counterwithin{figure}{section}
\counterwithin{table}{section}

\section{Appendices}
\label{app:appendices}
	This is to show how to handle appendices and figures in it (figure \ref{fig:Figure_app}).
\begin{figure}
	\centering
	\includegraphics[width=\linewidth]{fig_section} 
	\caption{This is to show appendix figures are numbered differently.}
	\label{fig:Figure_app}
\end{figure}
	
	Maecenas volutpat nulla at tellus commodo fermentum. Vestibulum non lacus et lectus consectetur blandit ac eget diam. Sed cursus eros vel mauris consectetur rhoncus non non velit. Vestibulum porta convallis mauris quis varius. Suspendisse eu odio non lacus imperdiet sollicitudin et ut orci. Cum sociis natoque penatibus et magnis dis parturient montes, nascetur ridiculus mus. Vivamus volutpat feugiat congue. Sed scelerisque nulla quis risus gravida mollis. Pellentesque suscipit varius est vitae eleifend. Nam euismod nisi nec diam gravida tristique. Pellentesque habitant morbi tristique senectus et netus et malesuada fames ac turpis egestas. Aenean convallis imperdiet ligula, sit amet finibus lectus maximus eget. Quisque non euismod ex, interdum dignissim tortor. Morbi facilisis metus dolor, eget rhoncus justo hendrerit sed.
	
	
\end{appendices}

\FloatBarrier

\bibliography{Ref_Example}

\end{document}